  
\documentclass[12pt]{report}

\textheight=22.6cm
\textwidth=17cm 
\voffset-1.5cm
\hoffset-1.5cm

\setlength{\parskip}{4mm}
\usepackage{color}
\usepackage[brazil]{babel}
\usepackage[utf8]{inputenc}
\usepackage{amsfonts} 
\usepackage{graphicx}
\usepackage{fancyhdr}
\usepackage{makeidx}
\usepackage{mathrsfs}
\usepackage{indentfirst}
\usepackage[Conny]{fncychap}
\usepackage{amsmath}
\usepackage{float}
\usepackage{subfiles}
\usepackage{pgfplots}
\usepackage{capt-of}
\pgfplotsset{compat=1.18, width=10cm} 

\newcommand{\ru}{\rule{3mm}{3mm}}
\newtheorem{definicao}{Definição}[section]
\newtheorem{exemplo}{Exemplo}[section]
\newtheorem{proposicao}{Proposição}[section]
\newtheorem{teorema}{Teorema}[section]
\newenvironment{prova}
{\par\noindent\textbf{Prova:}\ }
{\hfill$\square$\par}


\begin{document} 
	
	\thispagestyle{empty}
	\begin{center}UNIVERSIDADE ESTADUAL PAULISTA JÚLIO DE MESQUITA
		FILHO\\
		Faculdade de Engenharia de Ilha Solteira\\
		Departamento de Matemática 
	\end{center}
		
	\vspace*{4cm} 
	\begin{center} 
		{\LARGE \textbf{Propriedades Dinâmicas e Topológicas do Mapa Logístico}}. 
	\end{center} 
	\vspace*{2cm}
	\begin{center}
		{\large Aluno: Pedro Henrique da Silva Zonta | pedro.zonta@unesp.br}
	\end{center} 
	\begin{center}
		{\large Orientador: Prof. Dr. Danilo Antonio Caprio | danilo.caprio@unesp.br }
	\end{center} 
	\vspace*{2cm}
	\begin{flushright}
		\begin{minipage}{12cm}
			Relatório Final do projeto de Iniciação Científica orientado
			pelo Prof. Dr. Danilo Antonio Caprio, realizado no Departamento de Matemática
			da  Universidade Estadual Paulista ``Júlio de Mesquita Filho'' - Faculdade de Engenharia
			de Ilha Solteira.
		\end{minipage}
	\end{flushright}
	\vspace*{3cm} 
	\begin{center} Ilha Solteira - SP, Agosto/2026 \end{center}
		
	\newpage
	\vspace*{0.75\textheight}
	\pagenumbering{arabic} 
	
	\chapter*{Plano de Trabalho:}
	
	\begin{center}
		\begin{tabular}{|c|l|} \hline
			& {\bf Uma Introdução a Sistemas Dinâmicos Caóticos} \\ \hline
			01 & Exemplos de sistemas dinâmicos \\ \hline
			02 & Definições elementares na teoria dos sistemas dinâmicos discretos \\ \hline
			03 & Hiperbolicidade \\ \hline
			04 & A família quadrática \\ \hline
			05 & Dinâmica simbólica \\ \hline
			06 & Conjugação topológica \\ \hline
			07 & Caos \\ \hline
		\end{tabular}
	\end{center}
	
	\newpage
	{\LARGE\textbf{Resumo}}
	
	Neste projeto, inicialmente estudaremos exemplos de sistemas dinâmicos para entendermos sua aplicabilidade e como um sistema desse tipo pode ser analisado. Essa análise será explanada posteriormente pelo estudo do comportamento assintótico da função que rege o sistema, que significa analisar o que acontece com o sistema com o passar do tempo.
	
	Posteriormente, investigaremos a dinâmica simbólica e o comportamento do mapa logístico no conjunto de Cantor $\Lambda$. Analisaremos a aplicação de itinerário $S: \Lambda \to \Sigma_2$, demonstrando a continuidade e a existência de um homeomorfismo entre o espaço do conjunto de Cantor e o espaço de sequências $\Sigma_2$.
	
	Em seguida, aprofundaremos a análise do conceito de caos segundo a definição de Devaney. Para facilitar esse estudo, utilizaremos ferramentas de conjugação e semi-conjugação topológica, permitindo relacionar sistemas com modelos mais viáveis de se estudar.
	
	{\LARGE\textbf{Introdução}}
	
	A análise inicial de um sistema dinâmico consiste em investigar o comportamento assintótico de sua função, o que permite prever e descrever o destino das órbitas iteradas e a estabilidade de seus pontos. Algumas funções demonstram que pequenas variações em constantes ou entradas podem gerar resultados de extrema imprevisibilidade.
	
	Nesse contexto, o este trabalho apoia-se na análise de sistemas dinâmicos discretos, com foco na família logística $F_{\mu}(x) = \mu x(1-x)$. O estudo demonstra que a alteração do parâmetro $\mu$ dita mudanças na estrutura e na dinâmica do sistema, alterando o surgimento, o desaparecimento e o comportamento de pontos fixos e periódicos, o que caracteriza o fenômeno das bifurcações. Ao longo dos capítulos, explora-se como as bifurcações alteram a estabilidade do sistema, classificando os pontos em atratores ou repulsores.
	
	Conforme a complexidade do sistema aumenta, especialmente quando os parâmetros de entrada levam pontos para fora do domínio $[0,1]$, a pesquisa foca no comportamento do mapa logístico restrito ao conjunto de Cantor $\Lambda$. Como alternativa para a inviabilidade de analisar algebricamente essas iterações de forma direta na função $F_\mu$, introduzimos a dinâmica simbólica estruturada sobre o espaço de sequências $\Sigma_2$. Essa abordagem permite modelar as órbitas do sistema por meio de sequências infinitas de 0s e 1s, estabelecendo um modelo exato, equivalente e mais viável para descrever a rica estrutura do mapa logístico.
	
	Para garantir a validade formal desse modelo, o projeto utiliza os conceitos de semi-conjugação e conjugação topológica. O trabalho demonstra a existência de um homeomorfismo entre o espaço do conjunto de Cantor $\Lambda$ e o espaço de sequências $\Sigma_2$ por meio da aplicação de itinerário e da utilização do mapa \textit{shift}. Essa equivalência topológica preserva propriedades entre os mapas conjugados e permite uma análise do caos. Apoiados na definição de Devaney, concluímos que o mapa logístico é caótico, o qual se baseia em três pilares essenciais: a dependência sensível às condições iniciais, a transitividade topológica e a densidade dos pontos periódicos.

	
	\subfile{capitulos/capitulo1SistemasDinamicosDiscretos.tex}
	
	\subfile{capitulos/capitulo2DefinicoesElementares.tex}
	
	\subfile{capitulos/capitulo3Hiperbolicidade.tex}
	
	\subfile{capitulos/capitulo4FamiliaLogistica.tex}
	
	\subfile{capitulos/capitulo5DinamicaSimbolica.tex}
	
	\subfile{capitulos/capitulo6ConjugacaoTopologica.tex}
	
	\subfile{capitulos/capitulo7Caos.tex}
	
	

\end{document}